\documentclass[11pt,a4paper]{article}
\usepackage[T1]{fontenc}
\usepackage[utf8]{inputenc}
\usepackage{lmodern,microtype}
\usepackage[margin=26mm]{geometry}
\usepackage{amsmath,amssymb,amsthm,mathtools}
\usepackage{booktabs,longtable,array,enumitem}
\usepackage{xcolor,fancyhdr}
\usepackage[hidelinks]{hyperref}
\hypersetup{pdftitle={Lucky-Car Reciprocity via Prescribed Forests, and Four Further Investigations},pdfauthor={KEIO}}
\pagestyle{fancy}
\fancyhf{}
\fancyhead[L]{\small T91--T95: a new attempt}
\fancyhead[R]{\small KEIO}
\fancyfoot[C]{\thepage}
\setlength{\headheight}{14pt}
\setlength{\emergencystretch}{2em}
\allowdisplaybreaks
\newtheorem{theorem}{Theorem}[section]
\newtheorem{lemma}[theorem]{Lemma}
\newtheorem{proposition}[theorem]{Proposition}
\newtheorem{corollary}[theorem]{Corollary}
\theoremstyle{definition}
\newtheorem{definition}[theorem]{Definition}
\theoremstyle{remark}
\newtheorem{remark}[theorem]{Remark}
\newcommand{\F}{\mathbb F}
\newcommand{\Q}{\mathbb Q}
\newcommand{\Z}{\mathbb Z}
\newcommand{\C}{\mathbb C}
\DeclareMathOperator{\Aut}{Aut}
\DeclareMathOperator{\Cay}{Cay}
\begin{document}
\title{Lucky-Car Reciprocity via Prescribed Forests,\\
and Four Further Investigations}
\author{KEIO}
\date{T91--T95, new attempt\quad\textbar\quad27 October 2026}
\maketitle
\begin{abstract}
We construct a weight-preserving bijection between permutation gap choices
and parking-colored forests with a prescribed set of parent edges. Its
inverse is determined by a first hitting time of a color-prefix deficit.
Combined with an established containment-count formula, this proves the
full lucky-set reciprocity conjecture of Avila, Ferroni, and Morales,
without changing either parking protocol or its parameter range.
Four further investigations give a superpolylogarithmic lower bound for
the zero-residue Stern-polynomial counting function, exact Somos--8
coefficient identities in three families, minimum valencies of cubelike
graphical regular representations on two infinite subsequences of
dimensions, and a modulo-four derivation of Sun's adjacent Legendre-symbol
determinant formula. Complete statements, partial statements, source
discrepancies, and established methodological inputs are distinguished.
The manuscript contains proofs and reproducible exact certificates; it
does not certify five historically unprecedented principles or worldwide
priority for its new constructions.
\end{abstract}

\section*{Scope of this attempt}
This manuscript is a fresh investigation, not a renumbering of the earlier
T91--T95 manuscript. Its most substantial methodological contribution is
the explicit auxiliary bijection in T91 and its transfer to the unchanged
Conjecture~5.1 of Avila--Ferroni--Morales. The nearest-larger-element
framework has classical precedents, but the prescribed-parent,
color-preserving inverse must be proved; it is not supplied merely by
renaming a known theorem.

\begin{center}
\small
\renewcommand{\arraystretch}{1.22}
\begin{tabular}{@{}>{\raggedright\arraybackslash}p{.07\textwidth}>{\raggedright\arraybackslash}p{.30\textwidth}>{\raggedright\arraybackslash}p{.56\textwidth}@{}}
\toprule
Track & Original target & Result and exact limitation\\
\midrule
T91 & Avila--Ferroni--Morales, Conjecture 5.1 & Complete lucky-set identity for every positive parameter vector and every subset. Explicit auxiliary bijection and inverse. Historical priority remains a separate question.\\
T92 & Ulas, Question 2.10 & The zero-residue count exceeds every fixed power of the logarithm. The residue-one branch remains unresolved here.\\
T93 & Barry, Conjectures 11--13 & Complete bilinear identities for 11 and 12, and for the original continued fraction in 13. A conflicting printed generating function in 13 has a sign error. The mechanism is Hone's established theorem.\\
T94 & \'{A}rnad\'ottir--Guo, Problem 8.3 & Exact minimum valencies on two infinite subsequences, plus four smaller dimensions. The full minimum-valency function is not determined.\\
T95 & Sun, Conjecture 4.8 & The missing prime-congruence branch is derived for all four parameters; the other branch is already published. The needed signed inverse data are implicit in earlier explicit formulas.\\
\bottomrule
\end{tabular}
\end{center}

\paragraph{The strict success conditions.}
Each target above predates this work and is identified in a primary source.
T91 preserves the exact original conjecture and has an unconditional
all-parameter proof. A bounded search found no identical proof of its
prescribed-forest bijection or exact conjecture resolution. This supports
investigating novelty; it is not proof that no such literature exists.
No track is certified here as satisfying the further absolute demand for
a historically unprecedented foundational principle and worldwide first
resolution. In particular, partial answers in T92 and T94 are not counted
as complete closures, and the close prior-method overlap in T93 and T95
is disclosed explicitly.

Independent internal reviews found no gap in the included proofs. These
reviews are not external peer review or proof-assistant formalization.
Finite computations are used only in their stated roles. The polynomial
interpolation certificate in T93 is a finite exact proof with an explicit
degree bound; the other computational checks supplement the written
infinite arguments.



\section{T91: a Cartesian-gap bijection and lucky-car reciprocity}

The target in this retry is Conjecture~5.1 of Avila--Ferroni--Morales,
\emph{Luck and magic for Pitman--Stanley polytopes and parking functions},
arXiv:2603.19194v2 (8 April 2026)~\cite{RetryAFM2026}. It is distinct from the parking-polytope
magic-positivity conjecture attempted in the preceding round. We retain
the full original parameter range and the exact lucky-set statement.

Let $\mathbf y=(y_1,\ldots,y_n)\in\mathbb Z_{>0}^n$, put
$u_k=y_1+\cdots+y_k$, $u_0=0$, and $m=u_n$. Write
$\operatorname{PF}(\mathbf y)$ for the words $w\in[m]^n$ whose increasing
rearrangement satisfies $w_{(k)}\le u_k$. Two parking procedures act on
this same set of words. In the FHMS procedure the available sites are
$u_1,\ldots,u_n$, each with capacity one. In the AFM procedure the sites
$2,3,\ldots,m$ each have capacity one; site $1$ is unavailable. Cars arrive
in label order and take the first available site at least their preference.
A car is lucky precisely when its final site equals its preference; a
failed car is not lucky. Denote the resulting lucky sets by $L_F(w)$ and
$L_A(w)$, respectively. These are precisely the two procedures in the
source's Conjecture~5.1, including its $\ell=1$ convention.

For $J\subseteq[n]$ put $\overline J=\{n+1-j:j\in J\}$.

\begin{theorem}[The exact lucky-set identity]
For every positive integer vector $\mathbf y$ and every $S\subseteq[n]$,
\[
 \#\{w\in\operatorname{PF}(\mathbf y):L_F(w)=S\}
 =\#\{w\in\operatorname{PF}(\mathbf y):
               L_A(w)=[n]\setminus\overline S\}.
\]
Consequently the two one-variable lucky enumerators satisfy
\[
 \sum_w t^{|L_F(w)|}=\sum_w t^{n-|L_A(w)|}.
\]
\end{theorem}

We first solve a separate weighted enumeration problem. Its parameters
$y_1,\ldots,y_n$ can be indeterminates; thus the auxiliary statement is
stronger than its positive-integer evaluation needed below.

\subsection{The auxiliary Cartesian-gap problem}

Fix a set $J\subseteq[n]$ of vertices called prescribed vertices. For a
permutation $\tau$ of $[n]$, let $r_i$ be the position of $i$ and let $p_i$
be the largest position strictly less than $r_i$ occupied by a label
greater than $i$, with $p_i=0$ if there is none. Consider the following two
finite sets.

\begin{itemize}
\item $\mathcal A_J$ consists of $\tau$, together with an integer $c_i$
      satisfying $p_i<c_i\le r_i$ for every $i\notin J$.
\item $\mathcal B_J$ consists of an ordinary parking vector
      $c=(c_1,\ldots,c_n)$, meaning $c_{(i)}\le i$, together with a parent
      $P(i)$ for each $i\in J$. This parent is either a label $k>i$ with
      $c_k=c_i$, or a ghost root $*$, permitted only when $c_i=1$.
      The ghost is placed above every label. Vertices outside $J$ have no
      parent. Every parent chain therefore terminates at a vertex outside
      $J$ or at the ghost.
\end{itemize}

Both objects have weight $\prod_{i\notin J}y_{c_i}$.

\begin{lemma}[Cartesian gaps and parking-colored forests]
There is a weight-preserving bijection $\mathcal A_J\longleftrightarrow
\mathcal B_J$. In particular,
\begin{align}
 &\sum_{\tau\in\mathfrak S_n}
       \prod_{i\notin J}\left(\sum_{h=p_i+1}^{r_i}y_h\right)\label{eq:t91gap}\\
 &\qquad=
 \sum_{c\in\operatorname{PF}_n}
    \left(\prod_{i\notin J}y_{c_i}\right)
    \prod_{i\in J}
       \left(\#\{k>i:c_k=c_i\}+\mathbf1_{c_i=1}\right).\nonumber
\end{align}
\end{lemma}

\begin{proof}
For an object in $\mathcal A_J$, give each prescribed vertex $i$ the parent
occupying position $p_i$, or $*$ if $p_i=0$. Assign the ghost color $1$,
retain all unprescribed colors, and give each prescribed vertex its
parent's color, working in decreasing label order. Parent labels strictly
increase, so the procedure is well defined. For every vertex $i$ one has
$c_i\le r_i$: this is given for an unprescribed vertex, inherited from a
parent in an earlier position for a prescribed vertex, or follows from
$1\le r_i$ for a ghost child. Hence $c$ is an ordinary parking vector.
The map preserves all unprescribed colors and therefore the weight.

Here is a complete inverse, including the choice of the split position.
The construction applies to any finite ordered set of labels; its parking
colors are numbered from $1$ to the number $N$ of labels. The empty object
returns the empty permutation. For a nonempty object, let $v$ be the
largest label. Let $D$ consist of $v$ and all vertices whose parent chain
passes through $v$, and put $d=|D|$ and $a=c_v$. All vertices in $D$ have
color $a$, and every vertex in $D\setminus\{v\}$ is prescribed. If $v$
is prescribed, its parent must be $*$ and $a=1$.

Set $C(s)=\#\{i:c_i\le s\}$ and $H(s)=C(s)-s$. The parking condition is
$H(s)\ge0$ for $1\le s\le N$. Since $D$ has color $a$,
$C(a)\ge C(a-1)+d\ge a-1+d$, with $C(0)=0$. Thus $H(a)\ge d-1$.
Also $H(N)=0$, and $H(s+1)-H(s)\ge-1$. Consequently the integer
\[
 k=\min\{s\ge a:H(s)=d-1\}
\]
exists; before this first hit one has $H(s)\ge d$ for $a\le s<k$.
Define
\[
 A=\{i\notin D:c_i\le k\},\qquad
 B=(D\setminus\{v\})\cup\{i:c_i>k\}.
\]
Then $|A|=k-1$ and $|B|=N-k$.

On $A$, retain the colors and parents. For $s<a$ its color-prefix count is
$C(s)\ge s$; for $a\le s<k$ it is $C(s)-d\ge s$. Thus this is a parking
vector of length $k-1$. In particular no vertex of $A$ actually has color
$k$, since already $k-1$ of its vertices have color at most $k-1$.
Its parent edges stay in $A$ or reach the ghost: an edge has equal colors
at its ends, and an edge entering $D$ would put its starting vertex in $D$.

On $B$, replace the colors by
\[
 c'_i=\begin{cases}1,&i\in D\setminus\{v\},\\c_i-k,&c_i>k,\end{cases}
\]
and replace each parent $v$ by the ghost. Other parents are retained.
These are again valid increasing equal-color parent edges. No original
ghost child outside $D$ lies in $B$, since it would have color $1\le k$.
For $1\le q\le N-k$ the new prefix count is
\[
 (d-1)+C(k+q)-C(k)=C(k+q)-k\ge q.
\]
So this is a parking vector of length $N-k$.

Recursively construct the permutations $\tau_A$ and $\tau_B$ for these
two smaller objects, and output
\[
                         \tau_A\,v\,\tau_B.
\]
For a prescribed vertex in $A$, its nearest larger label to the left is
unchanged. In $B$, a recursive ghost parent becomes $v$, the largest
label at the split; other nearest-larger-left parents stay unchanged.
For an unprescribed vertex in $B$, the allowed color interval is shifted
by $k$, recovering its original color. The root $v$, if unprescribed,
has $p_v=0$ and $1\le a\le k=r_v$; if prescribed, its ghost parent is
correct. Thus the reconstructed permutation maps to the given parents
and colors.

For uniqueness, suppose an input permutation maps to the given forest,
and let the largest label $v$ occupy position $k$. All of $D\setminus
\{v\}$ lies to the right of $v$, because parents lie strictly to the left.
No other parent edge can cross $v$: it would then have a larger label
$v$ between its endpoints, contrary to the nearest-larger-left rule.
Every component outside $D$ lies on one side of $v$; a component attached
to the ghost lies on the left. An unprescribed vertex on the left has
color at most $k-1$. An unprescribed vertex on the right has color greater
than $k$, because $v$ is a larger label to its left. Colors are constant
on parent components. It follows that the left and right sets are exactly
$A$ and $B$ above and that $H(k)=d-1$. The left parking inequalities give
$H(s)\ge d$ for every $a\le s<k$. Hence $k$ is exactly the first-hit
position used by the inverse. Induction proves uniqueness of the entire
permutation. The retained unprescribed colors also recover every gap
choice, completing the bijection.
\end{proof}

\paragraph{A worked inverse.}
Take $\tau=(3,1,4,2)$, prescribe $J=\{1,2\}$, and choose the unprescribed
colors $c_3=1$, $c_4=2$. The prescribed parents are $1\to3$ and $2\to4$;
color inheritance gives $c=(1,2,1,2)$ and weight $y_1y_2$. In the inverse,
$v=4$, $D=\{2,4\}$, $d=2$, and $a=2$. Since $C(2)-2=2$ and $C(3)-3=1$,
the first-hit position is $k=3$. The left set is $\{1,3\}$ and reconstructs
as $(3,1)$; the right singleton $\{2\}$ reconstructs as $(2)$. Their
concatenation around $4$ recovers $(3,1,4,2)$.

\subsection{Transfer to the unchanged conjecture}

\begin{proof}[Proof of the theorem]
Fix $J\subseteq[n]$ and count FHMS words for which $J\subseteq L_F(w)$;
cars outside $J$ are unrestricted. Fix their final outcome permutation
$\tau$, where position $r$ means the sparse site $u_r$. When car $i$
arrives, the sites assigned to larger labels in $\tau$ are precisely the sites
that will be filled later. Its nearest still-empty sparse site to the left is
$u_{p_i}$, or $0$ if $p_i=0$. Therefore its allowable preferences are
exactly
\[
                        u_{p_i}<w_i\le u_{r_i}.
\]
If $i\in J$, there is one allowed lucky preference, namely $u_{r_i}$.
Otherwise there are $\sum_{h=p_i+1}^{r_i}y_h$ choices. These choices are
independent for a fixed outcome permutation: induction on arrival order
verifies that they produce that outcome. The left side of
\eqref{eq:t91gap} is consequently
\[
                      \#\{w:J\subseteq L_F(w)\}.
\]

We now use the already established Proposition~4.13 of
Avila--Ferroni--Morales, with $\ell=1$. In our notation its exact statement
is that, for any $K\subseteq[n]$,
\[
 \#\{w:L_A(w)\subseteq K\}
 =\sum_{q\in\operatorname{PF}_n}
     \prod_{j\in K}y_{q_j}
     \prod_{j\notin K}
        \left(\#\{h<j:q_h=q_j\}+\mathbf1_{q_j=1}\right).
\]
Take $K=[n]\setminus\overline J$ and reverse the indices of $q$.
Ordinary parking vectors are invariant under coordinate permutations, so
the resulting sum is the right side of \eqref{eq:t91gap}. We obtain, for
every $J$,
\[
 \#\{w:J\subseteq L_F(w)\}
 =\#\{w:L_A(w)\cap\overline J=\varnothing\}.
\]
These are all upper-set sums of the two candidate exact-set
distributions. Boolean-lattice M\"obius inversion therefore identifies
the distributions, giving precisely the displayed identity for each
$S\subseteq[n]$. Summing by $|S|$ gives the one-variable identity.
\end{proof}

\paragraph{Method and priority scope.}
The new part of this argument is the explicit weighted Cartesian-gap
bijection and its first-hit inverse. It supplies a separate auxiliary
enumeration problem and transfers its solution to the exact pre-existing
Conjecture~5.1. The proof uses the source's established containment-count
formula; it does not claim that formula or Boolean inversion as new.
Nearest-greater-element trees, parking functions, and first-return
decompositions are classical objects, so a claim that this is a
historically unprecedented foundational principle would require more
than the present proof. A bounded prior-literature search and independent
proof review must be reported separately from correctness.

In particular, Shin's \emph{A new bijection between forests and parking
functions}, arXiv:0810.0427v2, Section~4~\cite{RetryShin2008}, constructs a closest-larger-on-
the-right tree from ordinary parking outcomes and recovers labels by
descendant ranks determined by jumps. The construction here instead
specifies a subset of parent edges on the left, retains the colors at
every unprescribed vertex with independent weights $y_h$, and uses the
first-hit equation $C(k)-k=d-1$ to reconstruct the missing gap positions.
These are concrete distinctions between the constructions, not proof of
the absence of every possible earlier equivalent result.

The verifier \texttt{t91\_cartesian\_verify.py} exhaustively tests both
directions of the auxiliary bijection for every $J$ through $n=5$, and
compares the two original parking procedures on explicit positive
parameter vectors. These finite tests supplement the all-parameter proof.


\section{T92 retry: superpolylogarithmic growth in the zero-residue branch}

The original target is Question~2.10 of Ulas,
\emph{Strong arithmetic property of certain Stern polynomials}
(arXiv:1909.10844v1, 2019; published 2020)~\cite{retry:ulas}. Write
\[
 B_0(t)=0,\quad B_1(t)=1,\quad B_{2n}(t)=tB_n(t),\quad
 B_{2n+1}(t)=B_n(t)+B_{n+1}(t),
\]
and let $e(n)=\deg B_n$. The source defines
\[
 \Pi_{r,2}(x)=\#\{1\le n\le x:
 B_n(t)\equiv1+r\sum_{i=1}^{e(n)}t^i\pmod2\}.
\]
Its question asks whether $\Pi_{r,2}(x)/(\log_2x)^2\to\infty$ for
$r\in\{0,1\}$. We prove substantially more for the unchanged $r=0$
branch. The $r=1$ branch remains unresolved in this investigation.

\begin{theorem}[Zero-residue growth]
As $x\to\infty$,
\[
 \log\Pi_{0,2}(x)\ge
 \frac{(\log\log x)^2}{2\log2}
 -O(\log\log x\,\log\log\log x).
\]
In particular, for every fixed real $A>0$,
$\Pi_{0,2}(x)/(\log x)^A\to\infty$.
\end{theorem}

The auxiliary problem is to construct arbitrarily many independently
variable binary blocks whose Stern polynomial has every positive-degree
coefficient even. A recursively symmetric word solves this problem.
The exact matrix identity below is proved directly, not invoked as an
unverified closure principle.

\begin{lemma}[A characteristic-two symmetric product identity]
Work over $\mathbb F_2[t]$, and put $d=t+1$ and
\[
 F(u)=\begin{pmatrix}du&u\\u+d&du+1\end{pmatrix},\qquad
 N(q,\delta)=\begin{pmatrix}dq+\delta&q\\q+d(1+\delta)&dq+1\end{pmatrix}.
\]
Set $D=\det N(q,\delta)=t^2q^2+\delta$. Then
\[
 N(q,\delta)F(u)N(q,\delta)
   =N(q+uD,t^2q^2).
\]
Moreover, with $e=(d,1)$ and $v=(0,1)^{\mathsf T}$, one has
$eN(q,\delta)v=1$.
\end{lemma}
\begin{proof}
The determinant and the last identity follow immediately by multiplication.
For the product identity, its upper-right entry is $q+uD$, and its
upper-left entry is $d(q+uD)+t^2q^2$. The lower-right entry is
$d(q+uD)+1$, and the lower-left entry is
$q+uD+d(1+t^2q^2)$. These are the four entries on the right.
All equalities are polynomial identities in $\mathbb F_2[t,q,\delta,u]$.
\end{proof}

To read a binary digit from left to right, use
\[
 L=\begin{pmatrix}t&0\\1&1\end{pmatrix},\qquad
 R=\begin{pmatrix}1&1\\0&t\end{pmatrix}
\]
for digits zero and one, respectively. Beginning at $v$, the state after
the word for $n$ is $(B_n,B_{n+1})^{\mathsf T}$. Products are ordered in
the order of their action. For $a\ge1$, let
\[
 f_a=1^a01\,0^a.
\]
Its matrix is $L^aRLR^a=F(t^a)$ over $\mathbb F_2[t]$.
Indeed, $R^a=\left(\begin{smallmatrix}1&S_a\\0&t^a\end{smallmatrix}\right)$,
$L^a=\left(\begin{smallmatrix}t^a&0\\S_a&1\end{smallmatrix}\right)$,
where $S_a=1+t+\cdots+t^{a-1}$ and $dS_a=1+t^a$.

For positive integers $a_1,\ldots,a_j,c$, define words recursively by
\[
 W_0=\varnothing,\qquad W_i=W_{i-1}f_{a_i}W_{i-1}
 \quad(1\le i\le j),\qquad V_{j,\mathbf a,c}=W_j1^c01.
\]
The integer with binary expansion $V_{j,\mathbf a,c}$ is positive and odd.
The matrix of $W_0$ is $I=N(0,1)$. The lemma therefore shows inductively
that the matrix of every $W_j$ is some $N(q_j,\delta_j)$. The first row
of the matrix $RLR^c$ of the final word $1^c01$ equals $(d,1)$:
\[
 (1,0)RLR^c=(d,dS_c+t^c)=(d,1).
\]
Consequently
\[
 B_{[V_{j,\mathbf a,c}]_2}(t)\equiv1\pmod2.
\]
This proves the required polynomial congruence for every choice of the
positive integer parameters, without a degree restriction.

For clarity, the first new three-parameter instance, $j=2$, has the binary
run lengths
\[
 (a,1,1,a,b,1,1,b,a,1,1,a,c,1,1).
\]
The number of digits in this instance is $4a+2b+c+8$.
For general $j$, the number of digits is
\[
 \ell=2\sum_{i=1}^j2^{j-i}(a_i+1)+c+2.
\]
These parameterizations are injective. Each block $f_a$ has four runs,
starts in one and ends in zero; the final block has three runs and starts
in one. Adjacent blocks never merge. The number of runs of $V$ is exactly
$4(2^j-1)+3=2^{j+2}-1$, determining $j$. After separating the last three
runs, the parameters $a_i$ are recovered recursively from the central
four-run block of $W_i$ and either surrounding copy of $W_{i-1}$.

\begin{proof}[Proof of the theorem]
Put $L=\lfloor\log_2x\rfloor$. Every word of at most $L$ digits gives an
index below $2^L\le x$. On writing $a_i=1+b_i$ and $c=1+b_0$, the length
becomes
\[
 \ell=2^{j+2}-1+b_0+\sum_{i=1}^j2^{j-i+1}b_i,
 \qquad b_0,\ldots,b_j\ge0.
\]
Thus the exact generating function for the number of words of length
$\ell$ in the $j$th family is
\[
 \frac{z^{2^{j+2}-1}}{\prod_{i=0}^j(1-z^{2^i})}.
\]
For large $L$, choose
$j=\lfloor\log_2L-2\log_2\log L\rfloor$. Then $2^{j+2}-1\le L/2$.
For the variables having weights $2^i$, $0\le i\le j$, restrict them to
\[
 0\le b\le\left\lfloor\frac{L}{2(j+1)2^i}\right\rfloor.
\]
Their weighted sum is at most $L/2$. Since $\lfloor y\rfloor+1\ge y$,
there are at least
\[
 \frac{L^{j+1}}{[2(j+1)]^{j+1}2^{j(j+1)/2}}
\]
distinct indices. Taking logarithms and using
$j=\log L/\log2-O(\log\log L)$ gives
\[
 \log\Pi_{0,2}(x)\ge
 \frac{(\log L)^2}{2\log2}-O(\log L\log\log L).
\]
Replacing $L$ by $\lfloor\log_2x\rfloor$ proves the displayed bound.
It also implies the final assertion, since the quadratic term in
$\log\log x$ dominates $A\log\log x$ for every fixed $A$.
\end{proof}

\paragraph{Scope and novelty.}
This solves the $r=0$ instance of the original Question~2.10 with a stronger
quantitative bound; it does not solve the question for both residues.
Ulas's Theorem~2.3 gives a two-parameter family and a quadratic logarithmic
lower bound. The symmetric recursion here introduces arbitrarily many
independent parameters, which is the substantive extension. A bounded
search located no identical construction or growth result. Matrix
recurrences and characteristic-two cancellation are established ideas;
no historically unprecedented foundational principle is certified.


\section{T93: three exact Somos--8 families and their singular parameters}
\label{retry93:sec}

This section gives a computer-assisted proof of the Hankel recurrences in
Barry's Conjectures 11 and 12, and of the recurrence attached to the original
continued fraction in Conjecture 13 \cite{retry93:barry}. It does not claim a
new foundational principle: the all-index rank theorem used below is a
polynomial consequence of Hone's Theorem 5.5 \cite{retry93:hone}. The
coefficient identification and specialization argument are made explicit,
including parameters at which the published rational coefficients have
poles. A bounded search through 27 October 2026 did not locate an explicit
prior proof of these three precise coefficient formulas; that does not
certify priority.

\subsection{Exact targets and a source discrepancy}
For $j\in\{11,12,13\}$, let $g^{(j)}(x;r)$ be the formal series with
constant term one satisfying the respective \emph{original continued
fraction} in Barry's paper:
\begin{align*}
 g^{(11)}&=\frac{1}{1-x/(1-rx)-x^2-x^3g^{(11)}},\\
 g^{(12)}&=\frac{1}{1-x-x^2/(1-rx)-x^3g^{(12)}},\\
 g^{(13)}&=\frac{1}{1-x(1-(r-1)x)/(1-rx)-x^2-x^3g^{(13)}}.
\end{align*}
Write $g^{(j)}=\sum_{m\ge0}g_m^{(j)}x^m$ and
\[
 h_n^{(j)}(r)=\det(g_{a+b}^{(j)}(r))_{0\le a,b\le n}.
\]
All these entries and determinants are polynomials in $\mathbb Z[r]$.
Put $\mathcal A_j(x)=x^3a_j(x^{-1})$, where
\[
\begin{array}{c|l|l}
j&a_j(X)&b_j\\ \hline
11&X^3-(r+1)X^2-X+r&r+1\\
12&X^3-(r+1)X^2+(r-1)X&1\\
13&X^3-(r+1)X^2+(r-2)X+r&2
\end{array}
\]
Then each original continued fraction is exactly equivalent to
\begin{equation}\label{retry93:quadratic}
 \mathcal A_j(x)g^{(j)}(x)=1-rx+x^3(1-rx)g^{(j)}(x)^2.
\end{equation}

\paragraph{Conjecture 13 transcription.}
The displayed Catalan expression in the source has denominator
$1+(r+1)x+(r-2)x^2+rx^3$. Its original continued fraction instead gives
$1-(r+1)x+(r-2)x^2+rx^3$, as follows by multiplying its denominator by
$1-rx$. These are different formal series; for example their coefficients
of $x$ are respectively $-(2r+1)$ and $1$. We do not assert the false
identity between these two displays. The result below preserves the
original continued fraction and the printed Somos parameters. Thus it
settles that consistent recurrence assertion, while the entire collection
of displays labeled Conjecture 13 contains a separate sign error.

\subsection{The coefficients and the complete statement}
For $j=11$, define
\begin{align*}
 D_{11}={}&r^4-2r^3+8r^2+2r-9,\\
 A_{11}={}&r^8-8r^7+21r^6-40r^5+35r^4-24r^3+71r^2+8r,\\
 B_{11}={}&8(r^9-6r^8+17r^7-30r^6+15r^5-14r^4-r^3-14r^2),\\
 C_{11}={}&8(r-1)(r^{10}-2r^8+29r^7-32r^6+39r^5+18r^4+11r^3-r^2+r),\\
 E_{11}={}&2r^{13}-13r^{12}+48r^{11}-85r^{10}+83r^9-11r^8+124r^7\\
 &\quad-454r^6+364r^5-263r^4-84r^3-189r^2-25r-9.
\end{align*}
For $j=12$, define
\begin{align*}
 D_{12}={}&r^2-4r+3,\\
 A_{12}={}&r^4-11r^3+26r^2-16r-5,\\
 B_{12}={}&2r^5-19r^4+40r^3-13r^2-5r,\\
 C_{12}={}&(r-1)(3r^6-12r^5+15r^4+25r^3-62r^2-36r-5),\\
 E_{12}={}&r^9-8r^8+26r^7-43r^6+40r^5-17r^4-23r^3+27r^2+19r+3.
\end{align*}
For $j=13$, define
\begin{align*}
 D_{13}={}&2(r^3-3r^2-5r+7),\\
 A_{13}={}&-r^7+8r^6-25r^5+20r^4+37r^3-75r-28,\\
 B_{13}={}&(r+1)(r^8-11r^7+47r^6-83r^5+17r^4+71r^3+45r^2-169r+210),\\
 C_{13}={}&(r^2-1)(3r^8-29r^7+115r^6-225r^5+181r^4+105r^3\\
 &\quad-255r^2-235r+84),\\
 E_{13}={}&-r^{10}+17r^9-96r^8+212r^7-54r^6-594r^5+796r^4\\
 &\quad+36r^3-721r^2+329r+588.
\end{align*}
The rational tuple
$(A_j,B_j,C_j,E_j)/D_j$ is exactly the tuple in the respective original
conjecture; the common denominator form only multiplies numerator and
denominator of its third component by $r-1$ when $j=11,12$.

\begin{theorem}[Three polynomial Somos--8 identities]\label{retry93:target}
For each $j=11,12,13$, all $n\ge8$, and every $r\in\mathbb C$,
\begin{equation}\label{retry93:identity}
 D_j h_n^{(j)}h_{n-8}^{(j)}
 =A_jh_{n-1}^{(j)}h_{n-7}^{(j)}
 +B_jh_{n-2}^{(j)}h_{n-6}^{(j)}
 +C_jh_{n-3}^{(j)}h_{n-5}^{(j)}
 +E_j(h_{n-4}^{(j)})^2.
\end{equation}
Thus the printed rational-coefficient Somos--8 recurrence holds wherever
$D_j(r)\ne0$, in its division-free bilinear interpretation even when some
Hankel determinants vanish. At zeros of $D_j$, the polynomial identity
remains valid and gives the corresponding lower-order relation.
\end{theorem}

\subsection{A specialization lemma for exact recurrence identification}
The auxiliary question is how to identify coefficients of a recurrence
when one knows a universal rank bound, without assuming that every term
or every intermediate continued-fraction denominator is nonzero.
The following elementary answer is useful here but is not claimed as a
historically new principle.

\begin{lemma}\label{retry93:ranklemma}
Let $y_n,\mathbf v_n\in\mathbb Q[r]\times\mathbb Q[r]^4$, $n\ge8$.
Suppose every five-by-five minor of the matrix with rows
$(y_n,\mathbf v_n)$ vanishes identically. Let $M$ have rows
$\mathbf v_8,\ldots,\mathbf v_{11}$ and assume $\det M\not\equiv0$.
If $D\in\mathbb Q[r]$ and $\mathbf c\in\mathbb Q[r]^4$ satisfy
$\mathbf v_n\mathbf c=Dy_n$ for $8\le n\le11$, then this identity holds
for every $n\ge8$ and every specialization of $r$.
\end{lemma}
\begin{proof}
Work first over the field $\mathbb Q(r)$. The four columns of $M$ are
independent. The rank bound forces the first column of the full matrix
to be their unique linear combination, whose coefficients are determined
by the four displayed rows. Multiplication by $D$ gives the asserted
relation over $\mathbb Q(r)$. All its entries are polynomials, so it is a
polynomial identity and remains true under every specialization.
\end{proof}

\subsection{The all-index rank bound and its precise provenance}
We use the following consequence of Hone's Theorem 5.5 in the publisher
version (Theorem 5.4 in arXiv:1907.05204v3). Suppose $s_0=1$,
$s_1,s_2$ are arbitrary, and, for $m\ge3$,
\begin{align}\label{retry93:hone}
 s_m={}&u s_{m-2}+v s_{m-3}
 +w\sum_{i=0}^{m-2}s_i s_{m-2-i}
 +z\sum_{i=0}^{m-3}s_i s_{m-3-i}
 +t\sum_{i=0}^{m-4}s_i s_{m-4-i},
\end{align}
where an empty sum is zero. Put
$K_n=\det(s_{a+b})_{0\le a,b<n}$, $K_0=1$.
Then every five-by-five minor of the infinite matrix
\begin{equation}\label{retry93:rankmatrix}
 \left(K_nK_{n-8},\ K_{n-1}K_{n-7},\ K_{n-2}K_{n-6},\
 K_{n-3}K_{n-5},\ K_{n-4}^2\right)_{n\ge8}
\end{equation}
vanishes as a polynomial in $s_1,s_2,u,v,w,z,t$.
Indeed, Hone proves the existence of a common nonzero vector of Somos--8
coefficients for generic parameters in precisely this seven-parameter
moment recurrence. The entries of any fixed minor are polynomials in
those parameters. Its vanishing on the generic locus therefore proves
that it vanishes identically. This polynomial conclusion survives all
specializations, even ones outside the regular continued-fraction locus.
This use of an established theorem is the infinite-index part of the proof.

We now place the three original continued fractions within that polynomial
identity. Temporarily fix $j$ and abbreviate $a=a_j$, $b=b_j$, $g=g^{(j)}$.
Let $G(X)=X^{-1}g(X^{-1})$. Equation \eqref{retry93:quadratic} gives
\[
 G(X)=\frac{a(X)-Y(X)}{2(X-r)},\qquad
 Y^2=a(X)^2-4X(X-r)^2,
\]
with the formal branch for which $G=X^{-1}+X^{-2}+O(X^{-3})$.
Since $h_1=b$, its first normalized Jacobi tail is
\begin{equation}\label{retry93:tail}
 S(X)=\frac{X-1-G(X)^{-1}}{b}
      =\frac{P(X)-Y(X)}{2bX(X-r)},\qquad
 P(X)=2X(X-r)(X-1)-a(X).
\end{equation}
Over $\mathbb Q(r)$ the three values of $b$ are nonzero. The series $S$
has leading term $X^{-1}$. A direct rationalization yields
\begin{equation}\label{retry93:tailquadratic}
 P(X)S(X)-bX(X-r)S(X)^2
 =\frac{P(X)^2-Y(X)^2}{4bX(X-r)}=:R(X).
\end{equation}
The right-hand side is a polynomial of degree two: divisibility follows
from $P+a=2X(X-r)(X-1)$ and the displayed equation for $Y^2$.
The polynomial $P$ is monic cubic with coefficient $-(r+1)$ at $X^2$.
Consequently, after putting $X=Z+(r+1)/3$, write
\[
 P(X)=Z^3+pZ+q,\qquad X(X-r)=Z^2+\lambda Z+\mu,
 \qquad S(Z+(r+1)/3)=\sum_{m\ge0}s_mZ^{-m-1}.
\]
Comparing negative powers in \eqref{retry93:tailquadratic}, which have zero
coefficient on its right side, gives \eqref{retry93:hone} with
\[
 (u,v,w,z,t)=(-p,-q,b,b\lambda,b\mu).
\]
No geometric smoothness or term-by-term nonvanishing assumption is needed
for the polynomial rank identities just established.

For clarity, the determinant change caused by stripping this first tail
is exact. If $g(x)=1/(1-a_0x-bx^2s(x))$ with $g(0)=s(0)=1$, define the
Hankel kernel
\[
 \mathcal K_g(x,y)=\frac{xg(x)-yg(y)}{x-y}.
\]
The identity
\[
 \frac{\mathcal K_g(x,y)}{g(x)g(y)}=1+bxy\mathcal K_s(x,y)
\]
is a congruence by unit triangular coefficient matrices. Taking finite
principal determinants gives
\[
 \det(g_{a+c})_{0\le a,c\le n}
 =b^n\det(s_{a+c})_{0\le a,c<n}.
\]
Translation of the Cauchy variable is also a unit triangular binomial
change of the polynomial basis, so it preserves these Hankel determinants.
Therefore the determinants of the translated tail satisfy
\begin{equation}\label{retry93:gauge}
 h_n^{(j)}=b_j^nK_n.
\end{equation}
Every entry in row $n$ of \eqref{retry93:rankmatrix} is multiplied by the
same factor $b_j^{2n-8}$ under this change. Hence all five-by-five minors
for the corresponding rows formed from $h_n^{(j)}$ vanish over
$\mathbb Q(r)$, and then identically in $\mathbb Q[r]$.

\subsection{Finite exact coefficient certificate}
It remains to identify the four coefficients; this is a finite polynomial
calculation, rather than an inference from finitely many sequence terms.
Equation \eqref{retry93:quadratic} gives $\deg_r g_m^{(j)}\le m$ by induction.
The determinant expansion then gives
\[
 \deg_r h_n^{(j)}\le n(n+1).
\]
For each $j$, let $F_{j,n}(r)$ denote the difference between the two sides
of \eqref{retry93:identity}. Its degree is bounded by the maximum of
\[
 \deg D_j+n(n+1)+(n-8)(n-7)
\]
and the four numbers
\[
 \deg c_{j,k}+(n-k)(n-k+1)+(n-8+k)(n-7+k),\quad1\le k\le4,
\]
where $(c_{j,1},\ldots,c_{j,4})=(A_j,B_j,C_j,E_j)$.
The resulting bounds at $n=8,9,10,11$ are as follows.
\[
\begin{array}{c|rrrr|r}
j&8&9&10&11&\det M_j(2)\\\hline
11&76&96&120&148&-4757714578613175915968435866673495590560\\
12&74&94&118&146&9865855835586623025289744\\
13&75&95&119&147&-1346418219599972300425255472
\end{array}
\]
Here the four-by-four matrix $M_j$ has row $n$ equal to
$(h_{n-1}h_{n-7},h_{n-2}h_{n-6},h_{n-3}h_{n-5},h_{n-4}^2)$ for
$n=8,9,10,11$.

The accompanying standard-library program
\texttt{t93\_somos8\_verify.py} reconstructs the moments using
\eqref{retry93:quadratic}, evaluates all determinants by fraction-free
Bareiss elimination with every exact division asserted, and checks
\[
 F_{j,n}(0)=F_{j,n}(1)=\cdots=F_{j,n}(B_{j,n})=0,
\]
where $B_{j,n}$ is the tabulated degree bound. These are respectively
$B_{j,n}+1$ distinct integer arguments, so each $F_{j,n}$ is the zero
polynomial. The program also computes the three nonzero determinants
$\det M_j(2)$ shown above and independently checks them by rational
Gaussian elimination. All these exact checks pass; their machine-readable
outputs accompany the manuscript. This finite certificate proves the four
coefficient-identifying polynomial identities, with no numerical
approximation or conjectural extrapolation.

Apply Lemma~\ref{retry93:ranklemma} with
$y_n=h_n^{(j)}h_{n-8}^{(j)}$ and the four other products as
$\mathbf v_n$. The all-index rank bound has already been proved, the
nonzero witness shows $\det M_j\not\equiv0$, and the finite certificate
proves its four required identities. This proves
Theorem~\ref{retry93:target} for every $n$ and every $r$.

\subsection{Degeneracies and the boundary of the claimed novelty}
The denominator factorizations are
\[
 D_{11}=(r^2-1)(r^2-2r+9),\quad
 D_{12}=(r-1)(r-3),\quad
 D_{13}=2(r-1)(r^2-2r-7).
\]
In particular, the singular integer parameters give the following
Somos--6 relations, each valid in bilinear form for $m\ge7$:
\[
 h_mh_{m-6}=\alpha h_{m-1}h_{m-5}
             +\beta h_{m-2}h_{m-4}+\gamma h_{m-3}^2.
\]
\[
\begin{array}{c|r|rrr}
j&r&\alpha&\beta&\gamma\\\hline
11&-1&4&-8&8\\
11&1&4&0&8\\
12&1&1&0&5\\
12&3&-3&28&21\\
13&1&4&0&8
\end{array}
\]
These follow by substituting into the polynomial identity and setting
$m=n-1$. In particular, they remain valid when $h_1=0$ at $(j,r)=(11,-1)$;
no division by that determinant has been made in the final identity.
At the other algebraic zeros of the denominators, the same polynomial
identity gives an exact relation without evaluating an undefined rational
coefficient.

The new auxiliary task in this treatment was to identify a recurrence
uniformly through singular fibers from a rank bound and finite polynomial
certificates. Its solution is an application of elementary linear algebra,
Hankel congruence, interpolation, and the established genus-two theorem.
A direct proof of the universal rank bound from a new determinant
principle was not obtained. Consequently, although the exact recurrence
claims above are proved, this track does not meet the requested condition
of a historically new principle independent of those older mechanisms.


\section{T94: exact minimum valency on an infinite subsequence}

The unchanged target is Problem~8.3 of \'{A}rnad\'ottir and Guo,
\emph{Perfect State Transfer and Cayley Presentations},
arXiv:2608.20476v1 (20 August 2026)~\cite{retry:arna}. It asks for infinitude of cubelike
graphical regular representations with perfect state transfer and for their
minimum valency as a function of the dimension. Write $\mu(d)$ for this
minimum, when such a graph exists. The infinitude clause was addressed in
the previous round. This section obtains exact values of $\mu$ at an
explicit infinite sequence of dimensions, but does not determine $\mu(d)$
for arbitrary $d$. It is therefore a partial result, not a complete
resolution of the original problem.

\subsection{The auxiliary problem: erase a small rigid configuration}

For $m\ge1$, let $Q_m$ be the binary $m$-cube. If
$C\le\mathbb F_2^m$ is a binary linear code of minimum distance at least
three, its coset graph has vertices $\mathbb F_2^m/C$ and edges
$x+C\sim x+e_i+C$. This is a simple connected cubelike graph of valency
$m$. A coordinate permutation preserving $C$ is called a permutation
automorphism of the code.

\begin{lemma}[A necessary redundancy bound]
\label{t94:redundancy}
If a cubelike GRR on $2^d$ vertices has valency $m$, and $d\ge2$, then
\[
 m\le 2^{m-d}.
\]
\end{lemma}
\begin{proof}
Such a GRR must be connected. Indeed, if it has $c\ge2$ isomorphic
components of size $a$, then its automorphism group contains independent
regular actions on the components and all component permutations, and has
order at least $a^c c!>ac=2^d$ (including the edgeless case $a=1$, $c\ge4$).
Write its connection set as the distinct nonzero columns of a full-rank
$d\times m$ matrix $A$, and put $C=\ker A$. Then $C$ has dimension
$r=m-d$ and no word of weight one or two. Choose an $r\times m$ generator
matrix $B$ of $C$. If two columns of $B$ coincide, the transposition of
these coordinates fixes every word of $C$. It therefore induces an
automorphism of the coset graph fixing the zero coset. This automorphism
is nontrivial, because the two corresponding unit vectors have distinct
images modulo $C$. This contradicts the GRR condition. All $m$ columns
of $B$ are consequently distinct elements of $\mathbb F_2^r$, giving
$m\le2^r$.
\end{proof}

\begin{lemma}[Code rigidity and perfect state transfer]
\label{t94:quotient}
Let $C\le\mathbb F_2^m$ have minimum distance at least five. Then
\[
 \operatorname{Aut}(Q_m/C)
 = (\mathbb F_2^m/C)\rtimes\operatorname{Aut}_{\rm perm}(C).
\]
If the permutation group on the right is trivial and
$\boldsymbol1=(1,\ldots,1)\notin C$, the coset graph is a GRR with
perfect state transfer at time $\pi/2$.
\end{lemma}
\begin{proof}
For clarity, here is a direct proof of the standard quotient-cube
automorphism fact (also a special case of Fawcett's Theorem~1.4~\cite{retry:fawcett}).
At each vertex label its $m$ incident edges by their coordinate directions.
Every two-edge path with distinct directions lies in a unique four-cycle:
an additional four-cycle would give a nonzero codeword of weight at most
four. An automorphism $f$ induces a permutation $\pi_x$ of the directions
at each vertex $x$. Opposite edges in these unique squares show that
$\pi_x=\pi_{x+e_j}$ for every $j$. Connectedness makes this a single
permutation $\pi$. Following a path from zero yields
$f(x+C)=f(C)+\pi(x)+C$, and well-definedness gives $\pi(C)=C$.
Conversely, each such translation and coordinate permutation is an
automorphism. The action of the permutation group on the quotient is
faithful, since there are no codewords of weight two.

If $P_i$ denotes translation by $e_i+C$, these commuting involution
matrices satisfy
\[
 U(\pi/2)=\prod_{i=1}^m\exp(i\pi P_i/2)
 =i^mP_{\boldsymbol1+C}.
\]
This is perfect state transfer to a distinct vertex precisely when
$\boldsymbol1\notin C$.
\end{proof}

\begin{lemma}[Complementing a rigid set]
\label{t94:complement}
Suppose $R\subseteq\mathbb F_2^r\setminus\{0\}$ has trivial setwise
stabilizer in $\operatorname{GL}(r,2)$, and let $a=|R|$ satisfy
$2^{r-1}-a\ge5$ and $2^r-1-a>2^{r-1}$. Form a matrix $B$ whose columns
are exactly
\[
 Y=\mathbb F_2^r\setminus(R\cup\{0\}),\qquad m=2^r-1-a,
\]
and let $C$ be its row space. The same conclusion holds after adjoining a
single zero column and increasing $m$ by one. Then $Q_m/C$ is a cubelike GRR of valency
$m$ on $2^{m-r}$ vertices, with perfect state transfer at $\pi/2$.
\end{lemma}
\begin{proof}
For each nonzero row vector $u$, exactly $2^{r-1}$ vectors of
$\mathbb F_2^r$ satisfy $u\cdot y=1$. Consequently
\[
 \operatorname{wt}(uB)
 =2^{r-1}-|\{y\in R:u\cdot y=1\}|
 \in[2^{r-1}-a,2^{r-1}].
\]
In particular $B$ has rank $r$, its row code has minimum distance at
least five, and it does not contain the all-ones word, whose weight is
$m>2^{r-1}$. A coordinate permutation preserving the row space yields a
matrix $L\in\operatorname{GL}(r,2)$ with $LB=B\pi$. Thus $L$ preserves
the column set $Y$ and hence its nonzero complement $R$. It follows that
$L=I$. Since the columns of $B$ are distinct, $\pi$ is the identity. A single zero
column is fixed by every invertible linear map, so adjoining it changes none
of these conclusions or the nonzero codeword weights.
Apply Lemma~\ref{t94:quotient}.
\end{proof}

\subsection{An exact infinite sequence of minima}

\begin{theorem}
\label{t94:exact}
For every integer $r\ge10$ and $\varepsilon\in\{0,1\}$, put
\[
 d_{r,\varepsilon}=2^r-3r-3+\varepsilon.
\]
Then the minimum in the original Problem~8.3 satisfies
\[
 \boxed{\mu(d_{r,\varepsilon})=2^r-2r-3+\varepsilon
 =d_{r,\varepsilon}+r.}
\]
\end{theorem}
\begin{proof}
Construct a graph $\Gamma_r$ by subdividing the six edges
$01,02,03,12,13,23$ of $K_4$ into paths of respective lengths
\[
 2,\quad3,\quad1,\quad r-6,\quad1,\quad1.
\]
It has $r$ vertices and $r+2$ edges, is triangle-free, and has minimum
degree two. It is asymmetric: the four degree-three vertices and their
six connecting paths are intrinsic. The unique paths of lengths two and
three share endpoint $0$, so their endpoints fix $0,1,2$, and then $3$.
Each path is fixed pointwise. Label all vertices from $0$ to $r-1$ and set
\[
 R=\{e_i:0\le i<r\}
 \cup\{e_i+e_j:ij\in E(\Gamma_r)\}.
\]
Thus $|R|=2r+2$. A linear map preserving $R$ preserves the graph on $R$
in which $x$ and $y$ are adjacent when $x+y\in R$. A basis vertex $e_i$
has degree $2\deg_{\Gamma_r}(i)\ge4$ in this graph, while each edge
vertex $e_i+e_j$ has degree two. The triangle-free hypothesis rules out
adjacency between two edge vertices. Hence a linear map preserving $R$
permutes the basis and induces an automorphism of $\Gamma_r$. It fixes
every basis vertex and is therefore the identity.

The inequalities in Lemma~\ref{t94:complement} hold for every $r\ge10$.
That lemma, adjoining the zero column when $\varepsilon=1$, produces a
graph of valency $m_r=2^r-2r-3+\varepsilon$ and dimension
$m_r-r=d_{r,\varepsilon}$. Write $d_r=d_{r,\varepsilon}$ in the rest
of this proof. To prove optimality, suppose a competing GRR has valency
$m'<d_r+r$. Connectedness gives $m'\ge d_r$, while
$m'-d_r\le r-1$. Lemma~\ref{t94:redundancy} implies
\[
 m'\le2^{m'-d_r}\le2^{r-1}<d_r,
\]
a contradiction. The inequality $d_r>2^{r-1}$ holds at $r=10$ and
persists as $r$ increases.
\end{proof}

For $\varepsilon=0$, the first three pairs $(d_r,\mu(d_r))$ are
$(991,1001)$, $(2012,2023)$, and $(4057,4069)$.
Their graphs are specified by small binary generator matrices, without
enumerating their $2^{d_r}$ vertices.

\begin{proposition}[Four smaller exact values]
\label{t94:small}
One has
\[
 \mu(14)=19,\quad\mu(15)=20,\quad
 \mu(44)=50,\quad\mu(45)=51.
\]
\end{proposition}
\begin{proof}
Use the twelve-element connection set displayed immediately before
Problem~8.3 in the source, represented here by binary integers:
\[
 R_5=\{3,7,8,9,10,12,13,14,15,17,23,31\}\subset\mathbb F_2^5.
\]
Its linear setwise stabilizer is trivial. This also has the following
independent finite certificate, avoiding reliance on the published graph
automorphism calculation. In the Cayley graph with connection set $R_5$,
give zero color one and every other vertex color zero. Recursively replace
each color by the zero-based lexicographic rank of its current color
together with the sorted multiset of its neighbors' colors. The successive
numbers of color classes are $3,8,30,32$; the fourth colors, in vertex order
$0,1,\ldots,31$, are
\[
\begin{split}
 &(31,16,14,28,17,15,18,29,26,27,24,12,20,23,30,25,\\
 &\hspace{8pt}9,21,10,2,1,3,11,19,13,5,0,7,6,4,8,22).
\end{split}
\]
Every automorphism fixing zero preserves these colors and hence is the
identity; in particular the linear stabilizer of $R_5$ is trivial.

The code whose generator columns are the nineteen nonzero vectors
outside $R_5$ has nonzero weights ranging from $7$ to $13$. This is an
exact check of only $31$ row combinations, and the accompanying program
records the entire distribution. Therefore it has dimension five,
minimum distance at least five, and does not contain the all-ones word.
The proof of Lemma~\ref{t94:complement} applies (using these exact weights
in place of its coarse lower bound). This gives valency $19$ in dimension
$14$. Adjoining a zero column gives valency $20$ in dimension $15$.

For the next pair embed $R_5$ in the coordinate hyperplane
$\mathbb F_2^5\times\{0\}$ of $\mathbb F_2^6$, and add $e_5$ to obtain
$R_6$. Directly, every proper hyperplane of $\mathbb F_2^5$ contains at
most nine points of $R_5$. The old coordinate hyperplane is consequently
the unique hyperplane containing twelve points of $R_6$; every other
hyperplane contains at most ten. A linear map preserving $R_6$ must
preserve this hyperplane, fix its rigid set $R_5$, and fix the unique
outside point $e_5$. It is the identity. Lemma~\ref{t94:complement} gives
valencies $50,51$ in dimensions $44,45$, since the minimum codeword
weight is at least $32-13=19$.

Optimality follows from Lemma~\ref{t94:redundancy}. In dimensions
$14,15$, any smaller competing valency would have redundancy at most
four and valency at least $14$, while the only candidates of redundancy
four are respectively $18,19$, both exceeding $2^4$; lower redundancies
also violate the bound. In dimensions $44,45$, a smaller valency has
redundancy at most five and is at least $44>2^5$.
\end{proof}

\paragraph{Why this is not the full minimum-valency function.}
The redundancy bound alone is not always attained. For example, if
$d=2^r-r$ with $r\ge3$, it implies $m\ge d+r=2^r$.
Equality would make the generator columns of the kernel code equal to all
vectors of $\mathbb F_2^r$. Every nonidentity element of
$\operatorname{GL}(r,2)$ would then induce a nontrivial coordinate
permutation and hence a graph automorphism fixing zero. Equality is
impossible, so $\mu(2^r-r)\ge2^r+1$. Determining which column-set sizes
admit the necessary rigidity, and controlling all graph automorphisms
when the kernel has short words, remain unresolved steps.

\paragraph{Prior methods and strict-success status.}
The coset-graph viewpoint, the distance-five lifting theorem, and the
commuting-involution PST identity are established methods. The rigid set
$R$ is reused from the previous research round; the new step here is to
erase it from the full binary vector space and use the remaining vectors
as code-generator columns. The result improves a prior upper bound to an
exact minimum on an explicit infinite subsequence. A bounded search did
not locate this exact formula, but it does not establish worldwide
priority or a historically new foundational principle. Since the complete
function demanded in Problem~8.3 remains undetermined, this track does
not satisfy all of the requested success conditions.



\section{T95: sign recovery and Sun's adjacent Legendre determinant}
\label{sec:t95sun}

The target in this attempt is Conjecture~4.8 of Sun's 2024 paper
\cite{t95sun2024}, with exactly the original matrix, parameters, and prime
range. Its $p\equiv1\pmod4$ branch is already Theorem~1.1 of
Yang--Zhang \cite{t95yang2026}. We prove its $p\equiv3\pmod4$ branch
below. This is a new proof supplied in this investigation, but it uses
published inverse-square identities and classical determinant tools.
Consequently it is not certified as a historically new principle.

Write $\chi(a)=(a/p)$, with $\chi(0)=0$, and let $h(-p)$ be the class
number of $\mathbb Q(\sqrt{-p})$.

\begin{samepage}
\begin{theorem}[The original equation (4.20) of Sun]
\label{thm:t95sun48}
For every prime $p>3$ with $p\equiv3\pmod4$, put
$n=(p-1)/2$, $L=\chi(2)$ and $s=(-1)^{(h(-p)-1)/2}$. Then
\begin{align*}
&\det\bigl[x+\chi(j-k+1)+\chi(j)y+\chi(k)z+\chi(jk)w
                    \bigr]_{0\le j,k<n}\\
&\qquad=x-L(w+y-z+1)-sz+(1+s-L)(wx-yz).
\end{align*}
This is an identity in $\mathbb Z[x,y,z,w]$, with no parameter exclusions.
\end{theorem}
\end{samepage}

The auxiliary question is how to determine the signs of quantities whose
squares are known to be one, without redoing their cyclotomic evaluations.
For the two quantities needed here, reduction modulo four resolves the
signs. We give all the finite matrix algebra.

\subsection{The two published magnitude identities}
Set
\[
 C=[\chi(i-j)]_{0\le i,j\le n},\qquad
 W=C^{-1},\qquad u=W\mathbf1.
\]
The following established identities are the external input:
\begin{equation}\label{eq:t95old}
 \det C=1,\qquad W_{n0}^{2}=1,\qquad u_0^{2}=1.
\end{equation}
The determinant evaluation is Vsemirnov's theorem
\cite{t95vsemirnov}. The square identities are contained in
Yang--Yang--Zhang, Proposition~2.1 \cite{t95yyy2026}. Their convention is
$C_{ij}=\chi(j-i)$, the negative of ours; changing this sign preserves
the determinant (the order is even) and both squares. Their relation
$\gamma_0=-\gamma_n$ transfers the stated endpoint square to $u_0$.
We use only these unconditional identities, not their target conjecture.

\subsection{Modulo-four recovery of both signs}

Let $H=[\chi(i-j+1)]_{0\le i,j<n}$ and $D=\det H$.
Deleting row zero and column $n$ from $C$ leaves $H$, so the cofactor
formula, $n$ odd, and $\det C=1$ give
\begin{equation}\label{eq:t95corner}
 W_{n0}=-D.
\end{equation}
In particular $D\in\{1,-1\}$ by \eqref{eq:t95old}.

Let $S$ have ones on its superdiagonal and zeros elsewhere, and put
$A=\mathbf1\mathbf1^T-S$, of order $n$. Direct solution of $Av=y$ gives
\[
 v_0=\sum_{i=0}^{n-2}y_i+(2-n)y_{n-1},\qquad
 v_j=y_{n-1}-y_{j-1}\quad(1\le j<n),
\]
and $\det A=1$. For example, subtracting the last row of $A$ from every
earlier row gives $-1$ in columns $1,\ldots,n-1$, respectively; expansion
along column zero proves the determinant assertion.
All nonzero entries of $H$ are odd, and its zeros are precisely on the
superdiagonal. Thus $H=A+2B$ for an integral matrix $B$. Expansion of a
determinant modulo four yields
\[
 D\equiv1+2\operatorname{tr}(A^{-1}B)\pmod4.
\]
The displayed inverse shows
\begin{align*}
 \operatorname{tr}(A^{-1}B)
 &=\sum_{j=0}^{n-2}B_{j0}+(2-n)B_{n-1,0}
      +\sum_{j=1}^{n-1}B_{n-1,j}\\
 &=\sum_{k=1}^{n-1}(\chi(k)-1)
       +(2-n)\frac{\chi(n)-1}{2}.
\end{align*}
The omitted superdiagonal terms are zero. Since $n$ is odd,
\[
 D\equiv\chi(n)=\chi(-1/2)=-L\pmod4.
\]
The two possible integers $1,-1$ have different residues modulo four,
so
\begin{equation}\label{eq:t95signD}
 D=-L,\qquad W_{n0}=L.
\end{equation}

Next, let $C_0$ be obtained from $C$ by replacing column zero with ones.
Cramer's rule gives $u_0=\det C_0$. Compare $C_0$ with
$A_0=\mathbf1\mathbf1^T-\operatorname{diag}(0,1,\ldots,1)$, of order
$n+1$. Here $\det A_0=(-1)^n=-1$, and
\[
 (A_0^{-1}y)_0=(1-n)y_0+\sum_{j=1}^n y_j,\qquad
 (A_0^{-1}y)_j=y_0-y_j\quad(1\le j\le n).
\]
Write $C_0=A_0+2B_0$. Column zero and the diagonal of $B_0$ vanish,
so
\[
 \operatorname{tr}(A_0^{-1}B_0)
 =\sum_{j=1}^n(B_0)_{0j}
 =\sum_{j=1}^n\frac{-\chi(j)-1}{2}.
\]
Consequently, with $Q=\sum_{j=1}^n\chi(j)$,
\begin{equation}\label{eq:t95ucongruence}
 u_0\equiv n-1+Q\pmod4.
\end{equation}
Dirichlet's quadratic-excess class-number formula \cite{t95davenport} gives
$Q=(2-L)h(-p)$. The integer $Q$ is odd, so this also implies that
$h(-p)$ is odd. If $p\equiv3\pmod8$, then $n\equiv1\pmod4$ and
$L=-1$; if $p\equiv7\pmod8$, then $n\equiv3\pmod4$ and $L=1$.
In either case \eqref{eq:t95ucongruence} becomes
$u_0\equiv-h(-p)\equiv-s\pmod4$. Together with $u_0^2=1$ this proves
\begin{equation}\label{eq:t95signu}
 u_0=-s.
\end{equation}

\subsection{The exact two-dimensional determinant}
Put $\eta=(\chi(0),\ldots,\chi(n-1))^T$ and $v=H^{-1}\mathbf1$.
Let $R$ reverse the $n+1$ coordinates. Since $RCR=-C$,
$Ru=-u$, in particular $u_n=-u_0$. Also
$\mathbf1^Tu=0$ because $W$ is skew-symmetric.
Extend $v$ by a last zero, writing $\widehat v=(v^T,0)^T$. The last
$n$ coordinates of $C\widehat v$ are one, by the definition of $H$.
For some scalar $t$,
\[
 \widehat v=u+tWe_0.
\]
Its last coordinate and \eqref{eq:t95signD} give $t=u_0/L$.
Skew-symmetry implies $\mathbf1^TWe_0=-u_0$, whence
\begin{equation}\label{eq:t95inverseentries}
 \mathbf1^TH^{-1}\mathbf1=-\frac{u_0^2}{L}=-L.
\end{equation}
The first row of $C$ is
$-(\chi(0),\ldots,\chi(n))$. Taking coordinate zero of
$C\widehat v=\mathbf1+te_0$ yields
\[
 \eta^TH^{-1}\mathbf1=-1-t=-1+Ls.
\]
Column one of $H$ equals $\eta$, so
$H^{-1}\eta=e_1$. In particular,
$\mathbf1^TH^{-1}\eta=\eta^TH^{-1}\eta=1$.
Thus for $U=(\mathbf1,\eta)$ we have
\[
 U^TH^{-1}U=
 \begin{pmatrix}-L&1\\-1+Ls&1\end{pmatrix}.
\]
The matrix in Theorem~\ref{thm:t95sun48} is
$H+UMU^T$, where $M=\left(\begin{smallmatrix}x&z\\y&w\end{smallmatrix}\right)$.
The rank-two determinant lemma now gives
\begin{align*}
 \det(H+UMU^T)
 &=-L\det\left(I_2+M
          \begin{pmatrix}-L&1\\-1+Ls&1\end{pmatrix}\right)\\
 &=x-L(w+y-z+1)-sz+(1+s-L)(wx-yz).
\end{align*}
This proves the theorem for all parameters, since the calculation is a
polynomial identity and the fixed matrix $H$ is invertible.

\paragraph{The complementary published branch.}
For completeness, the already established $p\equiv1\pmod4$ formula
\cite[Theorem~1.1]{t95yang2026} is as follows. Let $h(p)$ and
$\varepsilon_p>1$ be the class number and fundamental unit of
$\mathbb Q(\sqrt p)$, and define
\[
 \varepsilon_p^{(2-L)h(p)}=a'_p+b'_p\sqrt p,
 \qquad 2a'_p,2b'_p\in\mathbb Z.
\]
With $n=(p-1)/2$, the same original determinant equals
\begin{align*}
 &(pb'_p-a'_p)((w+1)x-yz)+L(wx-(y+1)z)\\
 &\hspace{25mm}+L(b'_p-a'_p)\{w(1-x)+(y+1)(z+1)\}.
\end{align*}
This branch is quoted as a prior theorem, not presented as a result first
obtained in this investigation. Combined with
Theorem~\ref{thm:t95sun48}, it gives the original full Conjecture~4.8.
The two prime-congruence branches are not counted as two different problems.

\paragraph{Success-criterion assessment.}
The original problem and its complete four-parameter statement are
preserved. The proof is unconditional, using the published identities
\eqref{eq:t95old}. A bounded primary-source search did not locate an
explicit prior proof of equation~(4.20). This does not establish worldwide
priority. The prior-art overlap is stronger than the square-only formulation
might suggest: equations (3.2), (5.3), and (5.5) of
\cite{t95yyy2026} already encode both signed boundary data. In their
orientation, (5.3) simplifies to $d=L$; the product in (5.5) gives
$\gamma_n=\prod_{j=1}^n\chi(j)=-s$. Our orientation has
$W_{n0}=d$ and $u_0=\gamma_n$. Thus the signs above are an alternative
modulo-four derivation of information implicit in the earlier calculation,
not new boundary invariants. Inverse-data reduction already appears in
the cited papers, and reduction modulo four is a classical technique.
This argument therefore supplies a concrete proof of the target, but it does
\emph{not} certify the requested historically unprecedented method.
The auxiliary sign-recovery problem must not be promoted into an
independent previously open named problem.

\paragraph{Exact verification.}
The standard-library program \texttt{t95\_sun48\_verify.py} computes
class numbers independently by enumerating reduced primitive quadratic
forms, evaluates the two sign determinants, and tests the full original
matrix at nine integer parameter choices for every prime
$3<p<104$ with $p\equiv3\pmod4$. These finite checks support the
algebra; the all-prime theorem follows from the proof above.


\appendix
\section{Reproducibility and the status of the five investigations}

The accompanying archive contains the standalone source, the rendered
manuscript, five principal verification programs, their recorded exact
outputs, and an independent additional check of T91. All programs use
Python's standard library except the symbolic identity check in T92,
which uses SymPy~1.14.0. Run Python without the optimization flag
\texttt{-O}, so that the exact assertions remain active.

\paragraph{T91.}
\path{t91_cartesian_verify.py} checks both directions of the auxiliary
bijection for every prescribed subset through $n=5$. It also tests 400
seeded forest instances in larger dimensions, compares the two original
parking procedures on 23 positive parameter vectors, and includes
negative controls that fail when the left-parent or first-hit rule is
changed. The separate \path{t91_independent_random_audit.py} uses a
monotone-stack implementation and 720 cases with nonconsecutive labels
through $n=80$. The all-parameter result follows from the written inverse
construction, not the number of checked examples.

\paragraph{T92.}
\path{t92_stern_growth_verify.py} checks all four entries of the matrix
identity symbolically in characteristic two, independently computes 367
binary-word examples, and checks five integer Stern-polynomial
recurrences. The asymptotic bound follows from injectivity and the exact
word-length count in the text. These checks do not settle residue one.

\paragraph{T93.}
\path{t93_somos8_verify.py} performs 1,320 exact evaluations with proved
degree bounds, thereby certifying twelve coefficient-identifying
polynomial identities. It also verifies three nonzero rank witnesses and
compares independent moment and determinant implementations. Hone's
theorem supplies the universal rank identity; the rank-specialization
lemma then proves the conclusion for all indices. This is a
computer-assisted exact proof, not extrapolation from initial terms.

\paragraph{T94.}
\path{t94_cubelike_verify.py} checks the sparse rigid construction on
55 parameter values, the complete codeword weights in small instances,
and the finite rigidity certificates used for the four smaller exact
minima. The general subgroup, code-distance, and redundancy arguments
are proved in the text. No graph with $2^{991}$ vertices is enumerated.

\paragraph{T95.}
\path{t95_sun48_verify.py} independently enumerates reduced primitive
quadratic forms for 13 class numbers and evaluates 117 instances of the
original four-parameter determinant. The proof for every prime uses the
specified published inverse identities and exact matrix algebra.

\paragraph{Prior searches and what is still open.}
The literature comparison through 27 October 2026 inspected the original
target papers, their cited mechanisms, related author manuscripts, and
indexed exact-title searches. No identical T91 prescribed-forest
construction or explicit proof of its exact target was located. The
comparison with Shin is recorded in T91; T93 and T95 document more direct
overlap with earlier methods and formulas. These are bounded evidence
statements, not guarantees of worldwide priority.

Hill--Luo--Trinh--Vindas-Mel\'endez Conjecture~8.1 on magic positivity of
all generalized parking-function polytopes was also examined. Its
all-dimensional statement was not proved or disproved in this attempt.
T91's target is a different pre-existing conjecture, identified openly at
the start of its section; the lucky-set result is not advertised as a
solution of Hill's conjecture. Likewise, the residue-one branch in T92
and the full minimum-valency function in T94 remain unresolved here.

The strict requested combination---a complete original problem solution,
a historically unprecedented foundational principle, and verified
worldwide priority---has not been certified. T91 supplies a concrete
complete proof with a substantive auxiliary construction suitable for
further independent mathematical and bibliographic review.

\begin{thebibliography}{99}
\raggedright


\bibitem{RetryAFM2026}
N.~Avila, L.~Ferroni, and A.~H.~Morales,
\emph{Luck and magic for Pitman--Stanley polytopes and parking functions},
arXiv:2603.19194v2, 8 April 2026, Proposition~4.13 and Conjecture~5.1.
\url{https://arxiv.org/html/2603.19194v2}.

\bibitem{RetryFHMS2025}
M.~Ferreri, P.~E.~Harris, L.~Martinez, and E.~Swartz,
\emph{Enumerating Vector Parking Functions and their Outcomes Based on
Specified Lucky Cars}, arXiv:2508.13917v2, 10 September 2025.
\url{https://arxiv.org/html/2508.13917v2}.

\bibitem{RetryShin2008}
H.~Shin, \emph{A new bijection between forests and parking functions},
arXiv:0810.0427v2, 3 October 2008, Section~4.
\url{https://arxiv.org/pdf/0810.0427}.

\bibitem{RetryHill2026}
C.~Hill, A.~Luo, V.~Trinh, and A.~R.~Vindas-Mel\'endez,
\emph{Lattice slices, Ehrhart polynomials, and magic positivity of
generalized parking-function polytopes}, arXiv:2607.15503v1,
16 July 2026, Conjecture~8.1.
\url{https://arxiv.org/html/2607.15503v1}.


\bibitem{retry:ulas}
M.~Ulas, \emph{Strong arithmetic property of certain Stern polynomials},
arXiv:1909.10844v1, 24 September 2019, Question~2.10 and Theorem~2.3.
\url{https://arxiv.org/abs/1909.10844v1}.

\bibitem{retry:arna}
A.~S. \'{A}rnad\'ottir and K.~Guo,
\emph{Perfect State Transfer and Cayley Presentations},
arXiv:2608.20476v1, 20 August 2026, Problem~8.3.
\url{https://arxiv.org/html/2608.20476v1}.

\bibitem{retry:fawcett}
J.~B.~Fawcett, \emph{Locally triangular graphs and normal quotients of
the $n$-cube}, J.~Algebraic Combinatorics \textbf{44} (2016), 119--130,
Theorem~1.4. DOI:10.1007/s10801-015-0659-1.
\url{https://arxiv.org/abs/1507.01503}.


\bibitem{retry93:barry}
P.~Barry, \emph{Conjectures on Somos 4, 6 and 8 sequences using Riordan
arrays and the Catalan numbers}, arXiv:2211.12637v1, 22 November 2022,
Conjectures 11--13. \url{https://arxiv.org/abs/2211.12637}.

\bibitem{retry93:hone}
A.~N.~W. Hone, \emph{Continued fractions and Hankel determinants from
hyperelliptic curves}, Communications on Pure and Applied Mathematics
\textbf{74} (2021), 2310--2347, Theorem 5.5 and equation (4.22) in the publisher version
(the Somos--8 theorem is Theorem 5.4 in arXiv:1907.05204v3).
DOI:10.1002/cpa.21923; arXiv:1907.05204.
\url{https://kar.kent.ac.uk/79346/11/cpa.21923%20%281%29.pdf}.


\bibitem{t95sun2024}
Z.-W. Sun, \emph{Problems and results on determinants involving Legendre
symbols}, arXiv:2405.03626v8 (28 July 2024), Conjecture~4.8,
equations (4.19)--(4.20).
\url{https://arxiv.org/pdf/2405.03626v8}.

\bibitem{t95yang2026}
Y. Yang and Y. Zhang, \emph{Two determinant evaluations in Sun's
conjectures involving Legendre symbols}, arXiv:2605.19517,
Theorem~1.1.
\url{https://arxiv.org/abs/2605.19517}.

\bibitem{t95yyy2026}
Y. Yang, G. Yang, and Y. Zhang, \emph{A proof of a conjecture of Zhi-Wei
Sun on a truncated Legendre-symbol determinant}, arXiv:2606.22548v2
(8 September 2026), Proposition~2.1.
\url{https://arxiv.org/pdf/2606.22548v2}.

\bibitem{t95vsemirnov}
M. Vsemirnov, \emph{On the evaluation of R. Chapman's ``evil
determinant''}, Linear Algebra Appl. \textbf{436} (2012), 4101--4106.
\url{https://doi.org/10.1016/j.laa.2011.08.039}.

\bibitem{t95davenport}
H. Davenport, \emph{Multiplicative Number Theory}, second edition,
Springer, 1980, Sections~1 and~9 (Dirichlet's class-number formula).

\end{thebibliography}
\end{document}
